% Original: PDF strana 10, gornji slajd.
\slajdid{03-s019}
\begin{frame}[t,label=03-s019]{Koder prioriteta sa 16 ulaza}
\setlength{\parskip}{0pt}
Prethodne ideje primenjujemo na mrežu većeg kapaciteta: koder prioriteta sa \textbf{16 ulaza}.
\par\vspace{1mm}
\begin{columns}[T,onlytextwidth]
\begin{column}{.46\textwidth}\centering
% element: 03-s019:e01
\begingroup
\newcommand{\KPSimbol}[3]{%
\begin{scope}[shift={(#2,#3)}]
\draw[DEvod] (0,0) rectangle (26,20);
\node at (13,17) {KP};
\foreach \n/\y in {3/16.9,2/12.1,1/7.3,0/2.5}{
 \node[anchor=west,inner sep=1pt] at (0,\y) {I\n};
 \draw[DEvod] (-2,\y)--(0,\y);
 \coordinate (#1-i\n) at (-2,\y);
}
\foreach \n/\x in {1/9,0/18}{
 \node[anchor=south,inner sep=1pt] at (\x,0) {A\n};
 \draw[DEvod] (\x,0)--(\x,-2);
 \coordinate (#1-a\n) at (\x,-2);
}
\node[anchor=east,inner sep=1pt] at (26,12.1) {GS};
\draw[DEvod] (26,12.1)--(28,12.1);
\coordinate (#1-gs) at (28,12.1);
\end{scope}}

\begin{tikzpicture}[x=.70mm,y=.60mm,font=\fontsize{9}{11}\selectfont]
\foreach \k/\y in {3/69,2/46,1/23,0/0}{
 \KPSimbol{p\k}{0}{\y}
 \foreach \n in {3,2,1,0}{
  \pgfmathtruncatemacro{\ii}{4*\k+\n}
  \draw[DEvod] (p\k-i\n)--++(-2,0) node[left,inner sep=1pt] {I\ii};
 }
}
\KPSimbol{f}{43}{34.5}
\foreach \k/\x in {3/38,2/36,1/32,0/34}
 \draw[DEvod] (p\k-gs)--(\x,0 |- p\k-gs)|-(f-i\k);
\draw[DEvod] (f-gs)--++(3,0) node[right,inner sep=1pt] {GS};
\draw[DEvod] (f-a1)--++(0,-1) node[below,inner sep=1pt] {A3};
\draw[DEvod] (f-a0)--++(0,-1) node[below,inner sep=1pt] {A2};
\end{tikzpicture}
\endgroup

\end{column}
\begin{column}{.51\textwidth}
% element: 03-s019:e02
Ako osnovni koder ima aktivan ulaz, aktivan je i njegov \(GS\).
\textbf{Drugi nivo bira prioritetnu grupu} prema tim GS signalima
i daje njene indeksne bitove \(A_3,A_2\).
\par\vspace{2mm}
% element: 03-s019:e03
Preostaje izbor \(A_1,A_0\) iz \emph{iste, izabrane grupe}.
\textcolor{DEisticanje}{\textbf{Izlazi se ne smeju direktno spajati.}}
Ni logičko spajanje neuslovljenih kodova ne daje validnu informaciju.
\par\vspace{2mm}
% element: 03-s019:e04
Potrebni su \textbf{dodatni ulazi i izlazi} koji omogućavaju prepoznavanje
validnih \(A_1,A_0\) i njihovo logičko objedinjavanje.
\end{column}
\end{columns}
\end{frame}
